
% ============================================================
% UNIDAD 1. NÚMEROS COMPLEJOS
% PARTE II: GEOMETRÍA, FORMA POLAR, EULER, DE MOIVRE,
% RAÍCES, TFA Y LOGARITMO COMPLEJO
% Versión ampliada, formal y con ejemplos para guía docente
% ============================================================

% ============================================================
\section{Geometría de los números complejos}
% ============================================================

\subsection{El plano complejo}

Sea
\[
z=a+bi.
\]

Como $z$ queda determinado de manera única por el par ordenado $(a,b)$,
podemos representar a $z$ como el punto
\[
(a,b)\in\mathbb{R}^2.
\]

Esta representación da lugar al \emph{plano complejo} o \emph{plano de Argand}.

En él:
\begin{itemize}
    \item el eje horizontal se denomina \emph{eje real};
    \item el eje vertical se denomina \emph{eje imaginario}.
\end{itemize}

El número
\[
z=a+bi
\]
se identifica entonces con el punto
\[
(a,b)
\]
y también con el vector que va del origen al punto $(a,b)$.

Esta identificación permite interpretar geométricamente muchas operaciones
complejas.

\textbf{Observación docente.}
Conviene insistir en que la identificación
\[
\mathbb{C}\leftrightarrow\mathbb{R}^2
\]
es una identificación de conjuntos y, además, de espacios vectoriales reales,
pero el producto complejo no coincide con el producto componente a componente
de $\mathbb{R}^2$.

\paragraph{Ejemplo 1.}
\[
z=3+2i
\]
se representa por el punto
\[
(3,2).
\]

\paragraph{Ejemplo 2.}
\[
z=-4+i
\]
se representa por
\[
(-4,1).
\]

\paragraph{Ejemplo 3.}
\[
z=-2-5i
\]
se representa por
\[
(-2,-5).
\]

\paragraph{Ejemplo 4.}
El número real
\[
z=6
\]
se representa por el punto
\[
(6,0),
\]
sobre el eje real.

\paragraph{Ejemplo 5.}
El número imaginario puro
\[
z=-3i
\]
se representa por
\[
(0,-3),
\]
sobre el eje imaginario.

% ============================================================
\subsection{Módulo como distancia al origen}
% ============================================================

Sea
\[
z=a+bi.
\]

El módulo de $z$ es
\[
\boxed{|z|=\sqrt{a^2+b^2}.}
\]

Geométricamente, $|z|$ es la distancia euclidiana desde el origen $(0,0)$ al
punto $(a,b)$.

Por el Teorema de Pitágoras,
\[
|z|^2=a^2+b^2.
\]

Además, como
\[
z\overline z=a^2+b^2,
\]
se obtiene
\[
\boxed{|z|^2=z\overline z.}
\]

\textbf{Proposición.}
Para todo $z\in\mathbb{C}$:
\[
|z|\geq 0,
\]
y
\[
\boxed{|z|=0\iff z=0.}
\]

\paragraph{Ejemplo 1.}
\[
z=3+4i
\quad\Longrightarrow\quad
|z|=5.
\]

\paragraph{Ejemplo 2.}
\[
z=-5+12i
\quad\Longrightarrow\quad
|z|=13.
\]

\paragraph{Ejemplo 3.}
\[
z=7i
\quad\Longrightarrow\quad
|z|=7.
\]

\paragraph{Ejemplo 4.}
\[
z=-8
\quad\Longrightarrow\quad
|z|=8.
\]

\paragraph{Ejemplo 5.}
\[
z=1-\sqrt3\,i
\]
satisface
\[
|z|
=
\sqrt{1+3}
=
2.
\]

% ============================================================
\subsection{Distancia entre números complejos}
% ============================================================

Sean
\[
z_1=a+bi,
\qquad
z_2=c+di.
\]

La distancia entre ellos se define mediante
\[
\boxed{
d(z_1,z_2)=|z_1-z_2|.
}
\]

En efecto,
\[
z_1-z_2
=
(a-c)+(b-d)i,
\]
y por tanto
\[
|z_1-z_2|
=
\sqrt{(a-c)^2+(b-d)^2},
\]
que es exactamente la distancia euclidiana entre los puntos $(a,b)$ y $(c,d)$.

\paragraph{Ejemplo 1.}
Entre
\[
z_1=1+i,
\qquad
z_2=4+5i,
\]
la distancia es
\[
|z_1-z_2|
=
|-3-4i|
=
5.
\]

\paragraph{Ejemplo 2.}
Entre
\[
z_1=-2+3i,
\qquad
z_2=1-i,
\]
\[
d(z_1,z_2)
=
|-3+4i|
=
5.
\]

\paragraph{Ejemplo 3.}
La distancia del complejo $z$ al origen es
\[
|z-0|=|z|.
\]

\paragraph{Ejemplo 4.}
La circunferencia de centro $z_0$ y radio $r>0$ se expresa como
\[
\boxed{|z-z_0|=r.}
\]

\paragraph{Ejemplo 5.}
La región
\[
|z-(2-i)|<3
\]
representa el disco abierto de centro
\[
2-i
\]
y radio $3$.

% ============================================================
\subsection{Propiedades del módulo}
% ============================================================

\textbf{Teorema.}
Para cualesquiera $z,w\in\mathbb{C}$:
\[
\boxed{|zw|=|z||w|.}
\]

\textbf{Demostración.}
Usando
\[
|z|^2=z\overline z
\]
y
\[
\overline{zw}=\overline z\,\overline w,
\]
tenemos
\[
|zw|^2
=
zw\,\overline{zw}
=
zw\,\overline z\,\overline w
=
(z\overline z)(w\overline w)
=
|z|^2|w|^2.
\]
Como ambos lados son no negativos,
\[
|zw|=|z||w|.
\]
\hfill$\square$

Como consecuencia, si $w\neq0$,
\[
\boxed{
\left|\frac zw\right|
=
\frac{|z|}{|w|}.
}
\]

También:
\[
\boxed{|\overline z|=|z|.}
\]

\paragraph{Ejemplo 1.}
Si
\[
z=1+i,
\qquad
w=2-i,
\]
entonces
\[
|z|=\sqrt2,
\qquad
|w|=\sqrt5.
\]
Por tanto
\[
|zw|=\sqrt{10}.
\]

\paragraph{Ejemplo 2.}
Para
\[
z=3-4i,
\]
\[
|\overline z|
=
|3+4i|
=
5
=
|z|.
\]

\paragraph{Ejemplo 3.}
Si
\[
z=2+2i,
\qquad
w=1-i,
\]
entonces
\[
\left|\frac zw\right|
=
\frac{2\sqrt2}{\sqrt2}
=
2.
\]

\paragraph{Ejemplo 4.}
\[
|i^n|=1
\]
para todo entero $n$, porque $|i|=1$.

\paragraph{Ejemplo 5.}
Si
\[
|z|=3
\]
y
\[
|w|=4,
\]
entonces
\[
|z^2w^3|
=
|z|^2|w|^3
=
9\cdot64
=
576.
\]

% ============================================================
\subsection{Desigualdad triangular}
% ============================================================

\textbf{Teorema.}
Para todo $z,w\in\mathbb{C}$,
\[
\boxed{
|z+w|
\leq
|z|+|w|.
}
\]

\textbf{Demostración.}
Partimos de
\[
|z+w|^2
=
(z+w)(\overline z+\overline w).
\]
Entonces
\[
|z+w|^2
=
|z|^2+|w|^2+z\overline w+\overline z w.
\]
Pero
\[
z\overline w+\overline z w
=
2\operatorname{Re}(z\overline w)
\leq
2|z\overline w|
=
2|z||w|.
\]
Por tanto,
\[
|z+w|^2
\leq
|z|^2+2|z||w|+|w|^2
=
(|z|+|w|)^2.
\]
Como ambos lados son no negativos,
\[
|z+w|\leq |z|+|w|.
\]
\hfill$\square$

Una consecuencia importante es la desigualdad triangular inversa:
\[
\boxed{
\bigl||z|-|w|\bigr|
\leq
|z-w|.
}
\]

\textbf{Demostración.}
Por la desigualdad triangular,
\[
|z|
=
|(z-w)+w|
\leq
|z-w|+|w|.
\]
De aquí,
\[
|z|-|w|
\leq
|z-w|.
\]
Intercambiando $z$ y $w$,
\[
|w|-|z|
\leq
|z-w|.
\]
Por tanto,
\[
\bigl||z|-|w|\bigr|
\leq
|z-w|.
\]
\hfill$\square$

\paragraph{Ejemplo 1.}
Sea
\[
z=1+i,
\qquad
w=1-i.
\]
Entonces
\[
|z+w|
=
|2|
=
2,
\]
mientras que
\[
|z|+|w|
=
2\sqrt2.
\]
Así,
\[
2\leq2\sqrt2.
\]

\paragraph{Ejemplo 2.}
Si
\[
z=2,
\qquad
w=3,
\]
entonces hay igualdad:
\[
|z+w|=5=|z|+|w|.
\]

\paragraph{Ejemplo 3.}
Si
\[
z=2,
\qquad
w=-3,
\]
entonces
\[
|z+w|=1<5.
\]

\paragraph{Ejemplo 4.}
Para
\[
z=3+4i,
\qquad
w=1,
\]
la desigualdad triangular inversa da
\[
\bigl|5-1\bigr|
\leq
|2+4i|,
\]
es decir,
\[
4\leq\sqrt{20}.
\]

\paragraph{Ejemplo 5.}
Si
\[
|z|<2
\]
y
\[
|w|<3,
\]
entonces
\[
|z+w|
<
5.
\]

% ============================================================
\subsection{Argumento de un número complejo}
% ============================================================

Sea
\[
z=a+bi\neq0.
\]

Un \emph{argumento} de $z$ es cualquier ángulo $\theta$ que satisface
\[
a=|z|\cos\theta
\]
y
\[
b=|z|\sin\theta.
\]

Si $\theta$ es un argumento, entonces
\[
\theta+2k\pi,
\qquad
k\in\mathbb{Z},
\]
también lo es.

Por tanto,
\[
\boxed{
\arg z
=
\theta+2k\pi,
\qquad k\in\mathbb{Z}.
}
\]

El argumento es, por naturaleza, multivaluado.

\textbf{Definición.}
El \emph{argumento principal} de $z$, denotado por
\[
\operatorname{Arg}(z),
\]
es el único argumento que pertenece al intervalo
\[
(-\pi,\pi].
\]

Así,
\[
\boxed{
\arg z
=
\operatorname{Arg}(z)+2k\pi.
}
\]

\textbf{Observación docente.}
No debe utilizarse ciegamente
\[
\theta=\arctan\left(\frac ba\right),
\]
porque la función arctangente no distingue cuadrantes. Es necesario inspeccionar
los signos de $a$ y $b$, o utilizar una función de dos argumentos tipo
$\operatorname{atan2}(b,a)$.

\paragraph{Ejemplo 1.}
Para
\[
z=1+i,
\]
\[
\operatorname{Arg}(z)=\frac{\pi}{4}.
\]

\paragraph{Ejemplo 2.}
Para
\[
z=-1+i,
\]
el punto está en el segundo cuadrante y
\[
\operatorname{Arg}(z)=\frac{3\pi}{4}.
\]

\paragraph{Ejemplo 3.}
Para
\[
z=-1-i,
\]
\[
\operatorname{Arg}(z)=-\frac{3\pi}{4}.
\]

\paragraph{Ejemplo 4.}
Para
\[
z=-2,
\]
\[
\operatorname{Arg}(z)=\pi.
\]

\paragraph{Ejemplo 5.}
Para
\[
z=-3i,
\]
\[
\operatorname{Arg}(z)=-\frac{\pi}{2}.
\]

% ============================================================
\subsection{Forma polar o trigonométrica}
% ============================================================

Sea
\[
z=a+bi\neq0.
\]

Definimos
\[
r=|z|.
\]

Si $\theta$ es un argumento de $z$, entonces
\[
a=r\cos\theta,
\qquad
b=r\sin\theta.
\]

Sustituyendo en
\[
z=a+bi,
\]
obtenemos
\[
\boxed{
z
=
r(\cos\theta+i\sin\theta).
}
\]

Esta expresión se denomina \emph{forma polar} o \emph{forma trigonométrica}.

La conversión rectangular-polar se realiza mediante
\[
r=\sqrt{a^2+b^2},
\]
y una elección adecuada de $\theta$.

La conversión polar-rectangular se realiza mediante
\[
a=r\cos\theta,
\qquad
b=r\sin\theta.
\]

\paragraph{Ejemplo 1.}
\[
z=1+i.
\]
Entonces
\[
r=\sqrt2,
\qquad
\theta=\frac{\pi}{4}.
\]
Por tanto,
\[
z
=
\sqrt2
\left(
\cos\frac{\pi}{4}
+i\sin\frac{\pi}{4}
\right).
\]

\paragraph{Ejemplo 2.}
\[
z=-1+i.
\]
Entonces
\[
r=\sqrt2,
\qquad
\theta=\frac{3\pi}{4}.
\]
Así,
\[
z
=
\sqrt2
\left(
\cos\frac{3\pi}{4}
+i\sin\frac{3\pi}{4}
\right).
\]

\paragraph{Ejemplo 3.}
\[
z=-2-2i.
\]
Entonces
\[
r=2\sqrt2,
\qquad
\theta=-\frac{3\pi}{4}.
\]

\paragraph{Ejemplo 4.}
Convirtamos
\[
z
=
4\left(
\cos\frac{\pi}{3}
+i\sin\frac{\pi}{3}
\right)
\]
a forma rectangular:
\[
z
=
4\left(
\frac12+i\frac{\sqrt3}{2}
\right)
=
\boxed{2+2\sqrt3\,i}.
\]

\paragraph{Ejemplo 5.}
\[
z=5\left(\cos\pi+i\sin\pi\right)
=
-5.
\]

% ============================================================
\subsection{Producto y cociente en forma polar}
% ============================================================

Sean
\[
z_1=r_1(\cos\theta_1+i\sin\theta_1),
\]
\[
z_2=r_2(\cos\theta_2+i\sin\theta_2).
\]

Entonces
\[
\boxed{
z_1z_2
=
r_1r_2
\left[
\cos(\theta_1+\theta_2)
+i\sin(\theta_1+\theta_2)
\right].
}
\]

Por tanto:
\[
\boxed{
|z_1z_2|=|z_1||z_2|
}
\]
y
\[
\arg(z_1z_2)
=
\arg z_1+\arg z_2
\pmod{2\pi}.
\]

Si $z_2\neq0$,
\[
\boxed{
\frac{z_1}{z_2}
=
\frac{r_1}{r_2}
\left[
\cos(\theta_1-\theta_2)
+i\sin(\theta_1-\theta_2)
\right].
}
\]

\paragraph{Ejemplo 1.}
Sea
\[
z_1=2\left(\cos\frac{\pi}{6}+i\sin\frac{\pi}{6}\right),
\]
\[
z_2=3\left(\cos\frac{\pi}{4}+i\sin\frac{\pi}{4}\right).
\]
Entonces
\[
z_1z_2
=
6
\left(
\cos\frac{5\pi}{12}
+i\sin\frac{5\pi}{12}
\right).
\]

\paragraph{Ejemplo 2.}
Con los mismos números,
\[
\frac{z_1}{z_2}
=
\frac23
\left(
\cos\left(-\frac{\pi}{12}\right)
+i\sin\left(-\frac{\pi}{12}\right)
\right).
\]

\paragraph{Ejemplo 3.}
Si
\[
z_1=4e^{i\pi/3},
\qquad
z_2=2e^{i\pi/6},
\]
entonces
\[
z_1z_2
=
8e^{i\pi/2}
=
8i.
\]

\paragraph{Ejemplo 4.}
\[
\frac{6e^{i5\pi/6}}{3e^{i\pi/3}}
=
2e^{i\pi/2}
=
2i.
\]

\paragraph{Ejemplo 5.}
Si
\[
|z_1|=5,\quad \arg z_1=\frac{\pi}{7},
\]
\[
|z_2|=2,\quad \arg z_2=-\frac{\pi}{5},
\]
entonces
\[
|z_1z_2|=10
\]
y
\[
\arg(z_1z_2)
=
\frac{\pi}{7}-\frac{\pi}{5}
=
-\frac{2\pi}{35}
\pmod{2\pi}.
\]

% ============================================================
\subsection{Fórmula de Euler}
% ============================================================

La fórmula de Euler conecta la exponencial compleja con las funciones
trigonométricas.

Recordemos las series de Maclaurin:
\[
e^x
=
\sum_{n=0}^{\infty}\frac{x^n}{n!},
\]
\[
\cos x
=
\sum_{n=0}^{\infty}
(-1)^n\frac{x^{2n}}{(2n)!},
\]
\[
\sin x
=
\sum_{n=0}^{\infty}
(-1)^n\frac{x^{2n+1}}{(2n+1)!}.
\]

Sustituyendo $x=i\theta$:
\[
e^{i\theta}
=
1+i\theta+\frac{(i\theta)^2}{2!}
+\frac{(i\theta)^3}{3!}+\cdots.
\]

Separando términos reales e imaginarios:
\[
e^{i\theta}
=
\left(
1-\frac{\theta^2}{2!}
+\frac{\theta^4}{4!}
-\cdots
\right)
+
i
\left(
\theta-\frac{\theta^3}{3!}
+\frac{\theta^5}{5!}
-\cdots
\right).
\]

Por identificación con las series de $\cos\theta$ y $\sin\theta$,
\[
\boxed{
e^{i\theta}
=
\cos\theta+i\sin\theta.
}
\]

\textbf{Corolarios.}
\[
\boxed{
\cos\theta
=
\frac{e^{i\theta}+e^{-i\theta}}{2}
}
\]
y
\[
\boxed{
\sin\theta
=
\frac{e^{i\theta}-e^{-i\theta}}{2i}.
}
\]

\paragraph{Ejemplo 1.}
\[
e^{i0}=1.
\]

\paragraph{Ejemplo 2.}
\[
e^{i\pi/2}=i.
\]

\paragraph{Ejemplo 3.}
\[
e^{i\pi}=-1.
\]

\paragraph{Ejemplo 4.}
\[
e^{i3\pi/2}=-i.
\]

\paragraph{Ejemplo 5.}
\[
e^{i\pi/3}
=
\frac12+\frac{\sqrt3}{2}i.
\]

% ============================================================
\subsection{Forma exponencial}
% ============================================================

Usando Euler, la forma polar
\[
z=r(\cos\theta+i\sin\theta)
\]
puede escribirse como
\[
\boxed{
z=re^{i\theta}.
}
\]

Debido a la periodicidad
\[
e^{i(\theta+2k\pi)}
=
e^{i\theta},
\]
la representación exponencial completa es
\[
\boxed{
z=re^{i(\theta+2k\pi)},
\qquad
k\in\mathbb{Z}.
}
\]

\textbf{Proposición.}
Si
\[
z_1=r_1e^{i\theta_1},
\qquad
z_2=r_2e^{i\theta_2},
\]
entonces
\[
z_1z_2
=
r_1r_2e^{i(\theta_1+\theta_2)}.
\]

\paragraph{Ejemplo 1.}
\[
1+i
=
\sqrt2\,e^{i\pi/4}.
\]

\paragraph{Ejemplo 2.}
\[
-1+i
=
\sqrt2\,e^{i3\pi/4}.
\]

\paragraph{Ejemplo 3.}
\[
-2
=
2e^{i\pi}.
\]

\paragraph{Ejemplo 4.}
\[
-3i
=
3e^{-i\pi/2}.
\]

\paragraph{Ejemplo 5.}
\[
2e^{i\pi/6}
=
\sqrt3+i.
\]

% ============================================================
\subsection{Identidad de Euler}
% ============================================================

Tomando
\[
\theta=\pi
\]
en la fórmula de Euler:
\[
e^{i\pi}
=
\cos\pi+i\sin\pi
=
-1.
\]

Por tanto,
\[
\boxed{
e^{i\pi}+1=0.
}
\]

Esta identidad reúne las constantes
\[
0,\quad1,\quad e,\quad i,\quad\pi.
\]

\paragraph{Ejemplo 1.}
\[
e^{i\pi}=-1.
\]

\paragraph{Ejemplo 2.}
\[
e^{2\pi i}=1.
\]

\paragraph{Ejemplo 3.}
\[
e^{3\pi i}=-1.
\]

\paragraph{Ejemplo 4.}
\[
e^{(2k+1)\pi i}=-1.
\]

\paragraph{Ejemplo 5.}
\[
e^{2k\pi i}=1.
\]

% ============================================================
\subsection{Teorema de De Moivre}
% ============================================================

\textbf{Teorema de De Moivre.}
Para todo entero $n$,
\[
\boxed{
(\cos\theta+i\sin\theta)^n
=
\cos(n\theta)+i\sin(n\theta).
}
\]

Más generalmente, si
\[
z=r(\cos\theta+i\sin\theta),
\]
entonces
\[
\boxed{
z^n
=
r^n
\left(
\cos(n\theta)+i\sin(n\theta)
\right).
}
\]

\textbf{Demostración mediante Euler.}
Como
\[
z=re^{i\theta},
\]
entonces
\[
z^n
=
r^ne^{in\theta}.
\]
Aplicando Euler:
\[
e^{in\theta}
=
\cos(n\theta)+i\sin(n\theta).
\]
Por tanto,
\[
z^n
=
r^n
\left(
\cos(n\theta)+i\sin(n\theta)
\right).
\]
\hfill$\square$

\paragraph{Ejemplo 1.}
Calcular
\[
(1+i)^4.
\]
Como
\[
1+i=\sqrt2\,e^{i\pi/4},
\]
\[
(1+i)^4
=
(\sqrt2)^4e^{i\pi}
=
4(-1)
=
\boxed{-4}.
\]

\paragraph{Ejemplo 2.}
\[
\left(
\cos\frac{\pi}{6}
+i\sin\frac{\pi}{6}
\right)^3
=
\cos\frac{\pi}{2}
+i\sin\frac{\pi}{2}
=
i.
\]

\paragraph{Ejemplo 3.}
\[
(2e^{i\pi/3})^5
=
32e^{i5\pi/3}.
\]

\paragraph{Ejemplo 4.}
\[
(-1+i)^6.
\]
Como
\[
-1+i
=
\sqrt2\,e^{i3\pi/4},
\]
\[
(-1+i)^6
=
8e^{i9\pi/2}
=
8i.
\]

\paragraph{Ejemplo 5.}
Usando De Moivre:
\[
(\cos\theta+i\sin\theta)^2
=
\cos2\theta+i\sin2\theta.
\]
Comparando partes reales e imaginarias se recuperan:
\[
\cos2\theta=\cos^2\theta-\sin^2\theta,
\]
\[
\sin2\theta=2\sin\theta\cos\theta.
\]

% ============================================================
\subsection{Raíces $n$-ésimas de un número complejo}
% ============================================================

Queremos resolver
\[
w^n=z,
\]
donde
\[
z=re^{i\theta},
\qquad
r>0.
\]

Escribimos
\[
w=\rho e^{i\varphi}.
\]

Entonces
\[
w^n
=
\rho^ne^{in\varphi}.
\]

Para que $w^n=z$ debe cumplirse
\[
\rho^n=r
\]
y
\[
n\varphi
=
\theta+2k\pi.
\]

Por tanto,
\[
\rho=r^{1/n}
\]
y
\[
\varphi
=
\frac{\theta+2k\pi}{n}.
\]

Las $n$ raíces distintas son
\[
\boxed{
w_k
=
r^{1/n}
e^{i(\theta+2k\pi)/n},
\qquad
k=0,1,\ldots,n-1.
}
\]

Geométricamente, se encuentran sobre una circunferencia de radio
\[
r^{1/n}
\]
y están separadas por un ángulo constante
\[
\frac{2\pi}{n}.
\]

\paragraph{Ejemplo 1.}
Resolver
\[
w^2=1.
\]
Las raíces son
\[
w_0=1,
\qquad
w_1=-1.
\]

\paragraph{Ejemplo 2.}
Resolver
\[
w^2=-1.
\]
Como
\[
-1=e^{i\pi},
\]
\[
w_k
=
e^{i(\pi+2k\pi)/2}.
\]
Así,
\[
\boxed{w=\pm i}.
\]

\paragraph{Ejemplo 3.}
Resolver
\[
w^3=1.
\]
Las raíces son
\[
1,
\qquad
e^{2\pi i/3},
\qquad
e^{4\pi i/3}.
\]

\paragraph{Ejemplo 4.}
Resolver
\[
w^4=16.
\]
Como
\[
16=16e^{i0},
\]
las raíces tienen módulo
\[
16^{1/4}=2.
\]
Por tanto,
\[
w_k
=
2e^{ik\pi/2},
\qquad k=0,1,2,3.
\]
Así,
\[
\boxed{2,\ 2i,\ -2,\ -2i}.
\]

\paragraph{Ejemplo 5.}
Resolver
\[
w^3=8i.
\]
Tenemos
\[
8i=8e^{i\pi/2}.
\]
Entonces
\[
w_k
=
2e^{i(\pi/2+2k\pi)/3},
\qquad k=0,1,2.
\]

% ============================================================
\subsection{Raíces $n$-ésimas de la unidad}
% ============================================================

Consideremos
\[
z^n=1.
\]

Como
\[
1=e^{2k\pi i},
\]
las soluciones son
\[
\boxed{
\omega_k=e^{2\pi i k/n},
\qquad
k=0,\ldots,n-1.
}
\]

Estas raíces se distribuyen uniformemente sobre la circunferencia unidad y
forman los vértices de un polígono regular de $n$ lados.

Una raíz
\[
\omega=e^{2\pi i/n}
\]
se denomina raíz primitiva $n$-ésima de la unidad, y todas las raíces pueden
escribirse como
\[
1,\omega,\omega^2,\ldots,\omega^{n-1}.
\]

\textbf{Proposición.}
La suma de las raíces $n$-ésimas de la unidad, para $n>1$, es cero:
\[
\boxed{
1+\omega+\omega^2+\cdots+\omega^{n-1}=0.
}
\]

\textbf{Demostración.}
Como $\omega\neq1$ y $\omega^n=1$,
\[
1+\omega+\cdots+\omega^{n-1}
=
\frac{\omega^n-1}{\omega-1}
=
0.
\]
\hfill$\square$

\paragraph{Ejemplo 1.}
Para $n=2$:
\[
1,-1.
\]

\paragraph{Ejemplo 2.}
Para $n=3$:
\[
1,\quad
-\frac12+\frac{\sqrt3}{2}i,\quad
-\frac12-\frac{\sqrt3}{2}i.
\]

\paragraph{Ejemplo 3.}
Para $n=4$:
\[
1,\ i,\ -1,\ -i.
\]

\paragraph{Ejemplo 4.}
Las raíces quintas de la unidad son
\[
e^{2\pi ik/5},
\qquad
k=0,\ldots,4.
\]

\paragraph{Ejemplo 5.}
Para las raíces cúbicas de la unidad,
\[
1+\omega+\omega^2=0.
\]

% ============================================================
\subsection{Polinomios con raíces complejas}
% ============================================================

Los números complejos permiten estudiar de manera completa las raíces de
polinomios.

Si
\[
p(z)=a_nz^n+\cdots+a_1z+a_0,
\]
una raíz de $p$ es un número $z_0$ tal que
\[
p(z_0)=0.
\]

\textbf{Teorema del factor.}
\[
p(z_0)=0
\iff
(z-z_0)\text{ divide a }p(z).
\]

\paragraph{Ejemplo 1.}
\[
z^2+1=0
\]
tiene raíces
\[
\pm i.
\]

\paragraph{Ejemplo 2.}
\[
z^2-2z+5=0
\]
tiene raíces
\[
1\pm2i.
\]

\paragraph{Ejemplo 3.}
\[
z^3-1=0
\]
tiene las tres raíces cúbicas de la unidad.

\paragraph{Ejemplo 4.}
\[
z^4-16=0
\]
equivale a
\[
z^4=16,
\]
cuyas raíces son
\[
2,\ 2i,\ -2,\ -2i.
\]

\paragraph{Ejemplo 5.}
Si
\[
p(z)=z^2-6z+13,
\]
entonces
\[
z
=
\frac{6\pm\sqrt{36-52}}{2}
=
3\pm2i.
\]

% ============================================================
\subsection{Teorema Fundamental del Álgebra}
% ============================================================

\textbf{Teorema Fundamental del Álgebra.}

Todo polinomio no constante con coeficientes complejos posee al menos una raíz
en $\mathbb{C}$.

Equivalentemente, todo polinomio de grado $n\geq1$ puede factorizarse como
\[
\boxed{
p(z)
=
a_n
(z-z_1)\cdots(z-z_n),
}
\]
donde las raíces se cuentan con multiplicidad.

Por ello se dice que
\[
\boxed{
\mathbb{C}\text{ es algebraicamente cerrado}.
}
\]

\textbf{Observación docente.}
La demostración completa del Teorema Fundamental del Álgebra requiere
herramientas que exceden el nivel habitual de un primer curso de Álgebra
Lineal. Conviene enunciarlo, discutir su significado y usarlo, pero no
demostrarlo en esta unidad.

\paragraph{Ejemplo 1.}
Todo polinomio cuadrático complejo tiene dos raíces contando multiplicidad.

\paragraph{Ejemplo 2.}
\[
z^2+1
=
(z-i)(z+i).
\]

\paragraph{Ejemplo 3.}
\[
z^3-1
=
(z-1)(z-\omega)(z-\omega^2).
\]

\paragraph{Ejemplo 4.}
\[
z^2-2z+1
=
(z-1)^2.
\]
La raíz $z=1$ tiene multiplicidad $2$.

\paragraph{Ejemplo 5.}
Un polinomio de grado $5$ posee exactamente cinco raíces complejas contando
multiplicidad.

% ============================================================
\subsection{Raíces conjugadas de polinomios reales}
% ============================================================

\textbf{Teorema.}
Sea
\[
p(z)=a_nz^n+\cdots+a_0
\]
un polinomio con coeficientes reales. Si
\[
z_0
\]
es raíz de $p$, entonces
\[
\overline{z_0}
\]
también es raíz.

\textbf{Demostración.}
Si
\[
p(z_0)=0,
\]
entonces
\[
\overline{p(z_0)}=0.
\]
Como los coeficientes son reales,
\[
\overline{a_k}=a_k.
\]
Por tanto,
\[
\overline{p(z_0)}
=
p(\overline{z_0}),
\]
de donde
\[
p(\overline{z_0})=0.
\]
\hfill$\square$

\paragraph{Ejemplo 1.}
Si
\[
2+3i
\]
es raíz de un polinomio real, también lo es
\[
2-3i.
\]

\paragraph{Ejemplo 2.}
El polinomio
\[
z^2-4z+13
\]
tiene raíces
\[
2\pm3i.
\]

\paragraph{Ejemplo 3.}
El factor real asociado al par
\[
a\pm bi
\]
es
\[
(z-(a+bi))(z-(a-bi))
=
z^2-2az+(a^2+b^2).
\]

\paragraph{Ejemplo 4.}
Las raíces
\[
1\pm2i
\]
generan el factor
\[
z^2-2z+5.
\]

\paragraph{Ejemplo 5.}
Un polinomio real de grado impar tiene al menos una raíz real, porque las
raíces no reales aparecen en pares conjugados.

% ============================================================
\subsection{Exponencial compleja}
% ============================================================

Para
\[
z=x+iy,
\]
definimos
\[
e^z
=
e^x e^{iy}.
\]

Usando Euler:
\[
\boxed{
e^{x+iy}
=
e^x(\cos y+i\sin y).
}
\]

De aquí se obtiene
\[
\boxed{|e^z|=e^{\operatorname{Re}(z)}.}
\]

Además,
\[
e^{z+2\pi i}=e^z,
\]
por lo que la exponencial compleja es periódica con periodo $2\pi i$.

\paragraph{Ejemplo 1.}
\[
e^{1+i\pi}
=
e(-1)
=
-e.
\]

\paragraph{Ejemplo 2.}
\[
e^{2+i\pi/2}
=
e^2\left(\cos\frac{\pi}{2}+i\sin\frac{\pi}{2}\right)
=
ie^2.
\]

\paragraph{Ejemplo 3.}
\[
|e^{3-2i}|=e^3.
\]

\paragraph{Ejemplo 4.}
\[
e^{z+4\pi i}=e^z.
\]

\paragraph{Ejemplo 5.}
Resolver
\[
e^z=1.
\]
Las soluciones son
\[
\boxed{z=2k\pi i,\qquad k\in\mathbb{Z}.}
\]

% ============================================================
\subsection{Logaritmo complejo}
% ============================================================

Queremos resolver
\[
e^w=z,
\qquad
z\neq0.
\]

Escribimos
\[
z=re^{i\theta}
\]
y
\[
w=u+iv.
\]

Entonces
\[
e^w
=
e^u e^{iv}.
\]

Para que
\[
e^w=z,
\]
debe cumplirse
\[
e^u=r
\]
y
\[
v=\theta+2k\pi.
\]

Por tanto,
\[
u=\ln r
\]
y
\[
v=\theta+2k\pi.
\]

Así,
\[
\boxed{
\log z
=
\ln|z|
+
i(\theta+2k\pi),
\qquad
k\in\mathbb{Z}.
}
\]

El logaritmo complejo es multivaluado.

\textbf{Definición.}
El \emph{logaritmo principal} se define mediante
\[
\boxed{
\operatorname{Log}(z)
=
\ln|z|
+
i\operatorname{Arg}(z).
}
\]

\textbf{Observación docente.}
Conviene enfatizar que las identidades logarítmicas complejas requieren
cuidado por la multivaluación del argumento.

\paragraph{Ejemplo 1.}
Para
\[
z=1,
\]
\[
\log 1
=
2k\pi i.
\]

\paragraph{Ejemplo 2.}
Para
\[
z=-1,
\]
\[
\log(-1)
=
i(\pi+2k\pi).
\]
En particular,
\[
\operatorname{Log}(-1)=i\pi.
\]

\paragraph{Ejemplo 3.}
Para
\[
z=i,
\]
\[
\log i
=
i\left(
\frac{\pi}{2}+2k\pi
\right).
\]

\paragraph{Ejemplo 4.}
Para
\[
z=1+i,
\]
\[
|z|=\sqrt2,
\qquad
\operatorname{Arg}(z)=\frac{\pi}{4}.
\]
Entonces
\[
\operatorname{Log}(1+i)
=
\frac12\ln2
+
i\frac{\pi}{4}.
\]

\paragraph{Ejemplo 5.}
Para
\[
z=-2i,
\]
\[
|z|=2,
\qquad
\operatorname{Arg}(z)=-\frac{\pi}{2}.
\]
Por tanto,
\[
\operatorname{Log}(-2i)
=
\ln2-\frac{\pi}{2}i.
\]

% ============================================================
\subsection{Aplicaciones en ingeniería}
% ============================================================

Los números complejos aparecen de manera natural en numerosos modelos de
ingeniería.

Entre las aplicaciones más frecuentes se encuentran:
\begin{itemize}
    \item análisis de circuitos de corriente alterna;
    \item fasores;
    \item procesamiento de señales;
    \item transformadas de Fourier;
    \item sistemas de control;
    \item vibraciones;
    \item electromagnetismo;
    \item ecuaciones diferenciales;
    \item mecánica cuántica;
    \item comunicaciones.
\end{itemize}

La fórmula de Euler permite representar señales sinusoidales mediante
exponenciales complejas. Por ejemplo,
\[
A\cos(\omega t+\phi)
=
\operatorname{Re}
\left(
Ae^{i(\omega t+\phi)}
\right).
\]

Esta representación transforma muchas operaciones trigonométricas en
operaciones algebraicas más sencillas.

\paragraph{Ejemplo 1.}
Una señal
\[
v(t)=10\cos(\omega t)
\]
puede representarse mediante el fasor
\[
10e^{i0}.
\]

\paragraph{Ejemplo 2.}
\[
v(t)=5\cos(\omega t+\pi/3)
\]
corresponde al fasor
\[
5e^{i\pi/3}.
\]

\paragraph{Ejemplo 3.}
La multiplicación por
\[
e^{i\theta}
\]
representa geométricamente una rotación de ángulo $\theta$.

\paragraph{Ejemplo 4.}
Multiplicar por
\[
re^{i\theta}
\]
combina una dilatación de factor $r$ con una rotación de ángulo $\theta$.

\paragraph{Ejemplo 5.}
En análisis de señales, la identidad
\[
\cos(\omega t)
=
\frac{e^{i\omega t}+e^{-i\omega t}}{2}
\]
permite expresar una señal real como combinación de dos exponenciales
complejas.

% ============================================================
\subsection{Resumen conceptual de la segunda parte}
% ============================================================

La estructura geométrica y analítica de los números complejos puede resumirse
mediante
\[
\boxed{
\begin{array}{c}
z=a+bi
\\[0.2cm]
\downarrow
\\[0.2cm]
(a,b)\in\mathbb{R}^2
\\[0.2cm]
\downarrow
\\[0.2cm]
|z|,\ \arg z
\\[0.2cm]
\downarrow
\\[0.2cm]
z=r(\cos\theta+i\sin\theta)
\\[0.2cm]
\downarrow
\\[0.2cm]
z=re^{i\theta}
\\[0.2cm]
\downarrow
\\[0.2cm]
\text{De Moivre}
\\[0.2cm]
\downarrow
\\[0.2cm]
\text{potencias y raíces}
\\[0.2cm]
\downarrow
\\[0.2cm]
\text{raíces de la unidad}
\\[0.2cm]
\downarrow
\\[0.2cm]
\text{Teorema Fundamental del Álgebra}
\\[0.2cm]
\downarrow
\\[0.2cm]
e^z,\ \log z
\end{array}
}
\]

Esta segunda parte completa la transición entre la estructura algebraica de
$\mathbb{C}$ y su interpretación geométrica, preparando el terreno para
aplicaciones posteriores en Álgebra Lineal, ecuaciones diferenciales, señales
y sistemas de ingeniería.
